\documentclass[10pt,a4paper]{amsart}
\usepackage[margin=19mm]{geometry}
\usepackage{amsmath,amssymb,mathtools}
\usepackage[scr=rsfs,cal=euler]{mathalfa}
\usepackage{xfrac}
\usepackage[unicode,hidelinks,bookmarks=false]{hyperref}
\newcommand{\ie}{i.e.\ }
\newcommand{\cf}{cf.\ }
\newcommand{\resp}{respectively\ }
\newcommand{\Subsection}{\subsection}
\providecommand{\nobreakdash}{\nobreak\hbox{-}}
\setlength{\emergencystretch}{2em}
\multlinegap0pt
\usepackage{amssymb}
\usepackage{mathtools}
\usepackage{tikz}
\usepackage{mdframed}
\usepackage{tcolorbox}
\usepackage{enumerate}
\usepackage[shortlabels]{enumitem}
\usepackage{bm}
\newtheorem{proposition}{Proposition}
\newtheorem{lemma}{Lemma}
\usepackage{cleveref}
\crefname{proposition}{Proposition}{Propositions}
\crefname{lemma}{Lemma}{Lemmas}
\def\BL{\operatorname{BL}}
\newcommand\Lip{\operatorname{Lip}}
\newcommand\N{\mathbb{N}}
\newcommand\R{\mathbb{R}}
\newcommand\Z{\mathbb{Z}}
\DeclareMathOperator*{\argmin}{arg\,min}
\newcommand{\dist}[2]{{\operatorname{d}_{\operatorname{BL}}[#1;#2]}}
\renewcommand{\ge}{\geqslant}
\renewcommand{\le}{\leqslant}
\newcommand{\mb}{\mathbb}
\newcommand{\mbf}{\mathbf}
\title{The proportion of permutations fixing a $k$-set}
\author{Ben Green, Mehtaab Sawhney}
\date{Source: arXiv:2604.28116v1 (CC BY 4.0)}
\begin{document}
\maketitle

\noindent\emph{Source and licence.} The proportion of permutations fixing a $k$-set. Version arXiv:2604.28116v1 (CC BY 4.0), distributed by arXiv under the Creative Commons Attribution 4.0 International licence, http://creativecommons.org/licenses/by/4.0/, as recorded in the arXiv licence metadata for this exact version and shown on the abstract page https://arxiv.org/abs/2604.28116v1. Full licence notice: \texttt{licenses/arxiv-2604.28116-CC-BY-4.0.txt}.\par\smallskip

\noindent\textit{Editorial scope.} Counted unit: the article's proposition prop124 (source0.tex lines 2500--2502). The article's complete original proof of it is delivered with it (source0.tex lines 2503--2546, one \texttt{proof} environment of 44 lines). No mathematics is written, repaired or re-derived: the statement and the proof are the article's own lines, in the article's own notation, and no retained line is substituted or re-flowed.  Retained uncounted as the source-local context the proof needs: lines 921--921; lines 2444--2444; lines 2445--2445; lines 2446--2446; lines 2447--2447; lines 2496--2496; lines 3341--3343.  Nothing was omitted from any retained span, so the package carries no omission record.  No retained line contains a citation, so the wrapper carries no bibliography and no bibliography entry.  No retained line contains a figure environment, an image-inclusion command or drawing code, so no image is delivered and the package declares no asset dependency.  Editorial formatting only, in addition: an AMS \texttt{amsart} preamble that restates the article's own declarations and macros used by the retained lines, namely the environments \texttt{proposition}, \texttt{lemma} on separate counters, together with the standard packages those lines need.  The retained environments print in this document as Proposition 1, Lemma 1 in order of appearance, whereas the article prints the same items with the numbering of its own full document.  The complete native source of the article as distributed in the arXiv e-print for this exact version is bundled unchanged as \texttt{sources/199526/source0.tex}.
Suppose that $\beta$ is an upper walk which is $T$-positive up to $\ell$. Then we have \begin{equation}\label{lem73-first} \mb{E} \big| \varrho^*_{\mbf{u} \mid \beta}(\ell-1) - \varrho^*_{\mbf{u} \mid \beta}(\ell) \big|^2\ll 2^{3T - c\ell^{1/4}}.\end{equation}

\begin{equation}\label{exc-bd} \int_{\substack{\min \beta = -m \\ \argmin \beta \le \ell}} \mb{E} \varrho^*_{\mbf{u} \mid \beta}(\ell)^c\, \mathrm{d}\omega_m(\beta) \ll 2^{-m/4};\end{equation}

\begin{equation}\label{exc-rm-bd} \int_{\substack{\min \beta = -m \\ R_{\le N}(\beta) = r }} \mb{E} \varrho^*_{\mbf{u} \mid \beta}(\ell)^c\, \mathrm{d}\omega_m(\beta) \ll r^2(m + r + 1) e^{-r};\end{equation}

\begin{equation}\label{exc-qrm-bd} \int_{\substack{ \min \beta = -m \\ R_{\le N}(\beta) = r \\ \argmin \beta = q}} \mb{E} \varrho^*_{\mbf{u} \mid \beta}(\ell)^c\, \mathrm{d}\omega_m(\beta) \ll r^3 (1 + q)^{\kappa - 3/2} \end{equation} for $q \le \ell$, and

\begin{equation}\label{exc-mt-bd} \int_{\substack{\min \beta = -m \\ R_{\le N}(\beta) = r \\ T(\beta) > T}} \mb{E} \varrho^*_{\mbf{u} \mid \beta}(\ell)^c\, \mathrm{d}\omega_m(\beta) \ll r^{2}(1 + m)^{O(1)} T^{-\eta/2}.\end{equation}

\begin{equation}\label{mu-ell-def} \int \psi \, \mathrm{d} \mu^{(\ell)} := \sum_{m \ge 0}\int_{\substack{\min \beta = -m \\ \argmin \beta \le \ell}} \mb{E} \varrho^*_{\mbf{u} \mid \beta}(\ell)^c \psi (\log_2 \varrho^*_{\mbf{u} \mid \beta}(\ell)) \, \mathrm{d}\omega_m(\beta).\end{equation}

\begin{lemma}\label{completeness-measures}
$\mathcal{P}(\R/\Z)$ is complete with respect to the bounded Lipschitz distance.
\end{lemma}

\begin{proposition}\label{prop124}
There is a Borel measure $\mu$ on $\R/\Z$ with $\Vert \mu \Vert_{\BL} = \dist{\mu}{0} < \infty$ such that $\dist{\mu^{(\ell)}}{\mu} \rightarrow 0$ as $\ell \rightarrow \infty$.   
\end{proposition}

\begin{proof} Since the space of measures is complete with respect to the bounded Lipschitz metric (\cref{completeness-measures}) it is enough to show that the $\mu^{(\ell)}$ are a Cauchy sequence in the bounded Lipschitz norm. For the rest of the section suppose that $\psi \in C(\R/\Z)$ is a test function with $\Vert \psi \Vert_{\Lip} \le 1$.
We first observe the inequality 
\begin{equation}\label{lip-xy} \big| x^c \psi(\log_2 x) - y^c \psi(\log_2 y) \big| \ll |x - y|^c\end{equation} for $x,y \in [0,\infty)$. To prove this, suppose wlog that $y \ge x$ and write $y = x + h$. If $x = O(h)$ then the result is trivial. Otherwise, we have
\[ x^c \big( \psi (\log_2 x) - \psi(\log_2(x + h))\big) = x^c \big( \psi(\log_2 x) - \psi(\log_2 x + O(\frac{h}{x})) \big) \ll x^c \frac{h}{x} \ll h^c\] by the Lipschitz property. Also
\[ \big( (x + h)^c - x^c\big) \psi(\log_2(x + h)) \le \big| (x + h)^c - x^c\big| \ll x^c \big(1 + O(\frac{h}{x}) - 1\big) \ll h^c \big( \frac{h}{x}\big)^{1 - c} \ll h^c,\] and together these statements give \cref{lip-xy}. 
We note also that
\begin{equation}\label{simple-real-var} |x - y|^c \le x^c + y^c.\end{equation}

Let $\ell < \ell'$. Then by the definition \cref{mu-ell-def} and from \cref{lip-xy} we have
\begin{align}\nonumber\big| & \int \psi \, \mathrm{d}\mu^{(\ell)} - \int \psi \, \mathrm{d}\mu^{(\ell')} \big| \\\nonumber & \le \sum_m \int_{\substack{\min \beta = -m \\ \argmin \beta \le \ell}}\Big|  \mb{E} \varrho^*_{\mbf{u} \mid \beta}(\ell)^c \psi (\log_2 \varrho^*_{\mbf{u} \mid \beta}(\ell)) -  \mb{E} \varrho^*_{\mbf{u} \mid \beta}(\ell')^c \psi (\log_2 \varrho^*_{\mbf{u} \mid \beta}(\ell')) \Big| \, \mathrm{d}\omega_m(\beta) \\ \nonumber & \qquad\qquad\qquad\qquad\qquad\qquad + \sum_{m}\int_{\substack{\min \beta = -m \\ \ell < \argmin \beta \le \ell'}} \mb{E}\varrho^*_{\mbf{u} \mid \beta}(\ell')^c \, \mathrm{d}\omega_m(\beta)\\  
& \ll \sum_m \int_{\substack{\min \beta = -m \\ \argmin \beta \le \ell}}  \mb{E} \big|\varrho^*_{\mbf{u} \mid \beta}(\ell) -  \varrho^*_{\mbf{u} \mid \beta}(\ell')\big|^c \, \mathrm{d}\omega_m(\beta) + \sum_{m}\int_{\substack{\min \beta = -m \\ \ell < \argmin \beta \le \ell'}} \mb{E}\varrho^*_{\mbf{u} \mid \beta}(\ell')^c \, \mathrm{d}\omega_m(\beta).\label{key-integral}
\end{align}
To bound the first term we split into two parts. As usual, $R(\beta)$ denotes the least integer so that $\beta$ is $R(\beta)$-bounded, and $T(\beta)$ the least integer such that $\beta$ is $T(\beta)$-bounded. The `main' region of integration is over $m \le \ell^{\eta_0}$ and $\beta$ for which $R(\beta) \le \ell^{\eta_0}$ and $T(\beta) \le \ell^{1/8}$, where here $\eta_0\in (0,\eta)$ is an absolute constant to be specified. On this region, it follows from \cref{lem73-first} with a telescoping sum that $\mb{E} |\varrho^*_{\mbf{u} \mid \beta}(\ell) - \varrho^*_{\mbf{u} \mid \beta}(\ell')| \ll e^{-\Omega(\ell^{1/4})}$, and hence $\mb{E} |\varrho^*_{\mbf{u} \mid \beta}(\ell) - \varrho^*_{\mbf{u} \mid \beta}(\ell')|^c \ll e^{-\Omega(\ell^{1/4})}$ by H\"older's inequality. Since $\omega_m (\Z^{\N}_{\ge -m}) \ll m$, the contribution of the main region of integration to the first term in \cref{key-integral} is $\ll e^{-\Omega(\ell^{1/4})}$.

To estimate the contribution to the first term from the remaining portions, we first apply \cref{simple-real-var}, thereby reducing matters to estimating
\begin{equation}\label{remaining-portion-1} \sum_{m >\ell^{\eta_0}}\int_{\substack{\min \beta = -m \\\argmin \beta  \le \ell }}
\mb{E} \varrho^*_{\mbf{u} \mid \beta}(\tilde\ell)^c  \, \mathrm{d}\omega_m(\beta),\qquad \sum_{m \le \ell^{\eta_0}}\int_{\substack{\min \beta = -m \\
R(\beta) > \ell^{\eta_0}}}
\mb{E} \varrho^*_{\mbf{u} \mid \beta}(\tilde\ell)^c  \, \mathrm{d}\omega_m(\beta),\end{equation} 
and
\begin{equation}\label{remaining-portion-3} \sum_{m \le \ell^{\eta_0}}\int_{\substack{R(\beta) \le \ell^{\eta_0} \\ T(\beta) > \ell^{1/8}}}
\mb{E} \varrho^*_{\mbf{u} \mid \beta}(\tilde\ell)^c 1_{\min \beta = -m} \, \mathrm{d}\omega_m(\beta)\end{equation} for $\tilde\ell \in \{ \ell, \ell'\}$. (We dropped the $\argmin$ conditions in the second estimate in \cref{remaining-portion-1} and in \cref{remaining-portion-3} since we will not need these in the estimations.)

The contribution from the first integral in \cref{remaining-portion-1} is $\ll e^{-\Omega(\ell^{\eta_0})} < \ell^{-10}$ by \cref{exc-bd}.

The contribution from the second integral in \cref{remaining-portion-1} can be estimated using \cref{exc-rm-bd} with $N = \infty$, which gives a bound of 
\begin{equation}\label{mr-bd} \ll \sum_{m \le \ell^{\eta_0}} \sum_{r > \ell^{\eta_0}} r^2 (m + r + 1) e^{-r} \ll e^{-\Omega(\ell^{\eta_0/2})} < \ell^{-10}. \end{equation}

Finally, the contribution from \cref{remaining-portion-3} can be estimated using \cref{exc-mt-bd}, which gives a bound of
\[ \ll \ell^{-\eta/16} \sum_{m \le \ell^{\eta_0}} \sum_{r \le \ell^{\eta_0}} (r(1 + m))^{O(1)} \ll \ell^{-\eta/32},\]
provided that $\eta_0$ is chosen appropriately. This will be the weakest of the various estimates we obtain.

We also need to estimate the second term in \cref{key-integral}. Using \cref{exc-bd}, the contribution from $m > \ell^{\eta_0}$ is $< \ell^{-10}$ similarly to before.

The contribution from $m \le \ell^{\eta_0}$ and $r \ge \ell^{\eta_0}$ is $< \ell^{-10}$ as in the previous estimation of the second integral in \cref{remaining-portion-1}.

The final remaining contribution is at most
\[ \sum_{m \le \ell^{\eta_0}} \sum_{r \le \ell^{\eta_0}} \int_{\substack{\min \beta = -m \\ \ell < \argmin \beta \le \ell'}} \mb{E} \varrho^*_{\mbf{u} \mid \beta}(\ell')^c \, \mathrm{d}\omega_{m} (\beta).\] By \cref{exc-qrm-bd} with $N = \infty$ this is
\[ \ll \ell^{5\eta_0} \sum_{q > \ell} (1 + q)^{\kappa - 3/2} \ll \ell^{5\eta_0 + \kappa - 1/2} < \ell^{-1/10}.\]


Putting everything together, it follows that 
\[ \big| \int \psi \, \mathrm{d}\mu^{(\ell)} - \int \psi \, \mathrm{d}\mu^{(\ell')} \big| \ll \ell^{-\eta/32}.\] Since $\psi$ was an arbitrary Lipschitz test function, it follows that $\dist{\mu^{(\ell)}}{\mu^{(\ell')}} \ll \ell^{-\eta/32}$, which establishes that the $\mu^{(\ell)}$ are indeed a Cauchy sequence.
\end{proof}

\end{document}
